\begin{definition}[Known: finite-family Ramsey convention]\label{C:ramsey-definition}
We say that a small category $A$ has the finite-family Ramsey property if,
for every finite family $a_0,\ldots,a_k:p\to q$ and every integer $N\geq2$,
there is an object $r$ such that every map
$\chi:A(p,r)\to\{1,\ldots,N\}$ admits an arrow $b:q\to r$ with
$\chi(ba_0)=\cdots=\chi(ba_k)$.
The usual categorical Ramsey property makes all of $bA(p,q)$
monochromatic and therefore implies this condition.
\end{definition}

\begin{theorem}[Proved: a Ramsey sufficient condition for integral exactness]
\label{thm:ramsey-exact}
Let $A$ be a small category in which every finite family of objects has a
common successor. If $A$ has the finite-family Ramsey property, then $A$
satisfies condition III over $\Z$. If $A$ also satisfies condition I, then
$\colim_A:\Ab^A\to\Ab$ is exact.
\end{theorem}
\begin{proof}
Fix $a_0,\ldots,a_k:p\to q$. For $k=0$ the identity of $q$ suffices,
so we assume $k\geq1$. For every object $r$ we define a homomorphism
of abelian groups
\[
D_r:\Z[A(q,r)]\longrightarrow\Z[A(p,r)]^k,
\qquad D_r(t)=\bigl(t(a_i-a_0)\bigr)_{i=1}^k,
\]
and put $E=\sum_r\aug(\ker D_r)\subseteq\Z$.
We first verify that $E=\Z$ if and only if this family has an integral
affine equalizer. One direction is immediate. For the other, write
$1=\sum_{j=1}^m\aug(t_j)$ with $t_j\in\ker D_{r_j}$, absorbing any
integer multipliers into the $t_j$. Choose arrows $u_j:r_j\to s$ to a
common successor. Then $t=\sum_j u_jt_j\in\ker D_s$ and $\aug(t)=1$.

Suppose that there is no affine equalizer. We have $E=d\Z$ with either
$d=0$ or $d\geq2$. Let $\mathbb T=\mathbb R/\mathbb Z$ with its
translation-invariant circle metric. Choose $c=1/2\pmod\Z$ when $d=0$,
and $c=1/d\pmod\Z$ when $d\geq2$. Thus $c\ne0$ and $ec=0$ for
every $e\in E$. For every $r$ the formula
\[
\lambda_r(D_rt)=\aug(t)c
\]
defines a homomorphism $\lambda_r:\im D_r\to\mathbb T$: any two
preimages differ by an element of $\ker D_r$, whose augmentation kills
$c$.

We extend $\lambda_r$ to a homomorphism
$\widetilde\lambda_r:\Z[A(p,r)]^k\to\mathbb T$.
Here we use divisibility of the circle group. For completeness, a
homomorphism from a subgroup $B'$ of an abelian group $B$ to a divisible
group extends across one additional element $b\in B$: if no positive
multiple of $b$ lies in $B'$, choose its image arbitrarily; otherwise take
the least positive $n$ with $nb\in B'$ and choose an image $x$ with
$nx$ equal to the given image of $nb$. The extension to $B'+\Z b$ is
then well-defined. Zorn's lemma gives an extension to $B$.

Let $\phi_i^r:A(p,r)\to\mathbb T$ be the function obtained by restricting
$\widetilde\lambda_r$ to basis elements in its $i$th summand. For every
$b:q\to r$ we have
\begin{equation}\label{eq:circle-certificate}
\sum_{i=1}^k\left(\phi_i^r(ba_i)-\phi_i^r(ba_0)\right)=c.
\end{equation}
Put $\eta=\operatorname{dist}_{\mathbb T}(c,0)>0$ and choose
$\delta>0$ with $k\delta<\eta$. Partition $\mathbb T^k$ into finitely
many cells such that two points in one cell have distance less than
$\delta$ in every coordinate; we may arrange that the number $N$ of
cells is at least two. Choose $r$ from the finite-family Ramsey property
for this $N$. Color $h:p\to r$ by the cell containing
$(\phi_1^r(h),\ldots,\phi_k^r(h))$. There is $b:q\to r$ for which
$ba_0,\ldots,ba_k$ have one color. Every summand on the left of
\eqref{eq:circle-certificate} has circle distance less than $\delta$
from zero. Their sum has distance less than $k\delta<\eta$, contradicting
\eqref{eq:circle-certificate}. Therefore $E=\Z$, proving III.
The last assertion follows from \citet[Proposition~3.3]{IsbellMitchell1976}.
\end{proof}

\begin{corollary}[Proved: the Ramsey candidates in route C4 are exact]
\label{cor:ordered-graphs-exact}
The categories of nonempty finite linearly ordered simple graphs and of
nonempty finite linearly ordered triangle-free graphs, with order-preserving
induced embeddings, are exact over $\Z$ and hence over every commutative
ring. They are total and are not filtered. More generally, every nonempty
rigid Ramsey amalgamation class of finite structures with the joint
embedding property has exact abelian colimits.
\end{corollary}
\begin{proof}
We take a small skeleton in each case. Joint embedding gives common
successors, and amalgamation is condition I. The ordinary Ramsey property
for ordered graphs is \citet[Corollary~4.3]{HubickaNesetril2019}; for ordered
triangle-free graphs it is \citet[Theorem~15]{Hubicka2025}.
These statements concern all finite colorings and order-preserving induced
embeddings, exactly the morphisms used here. Because the structures are
rigid, copies and embeddings have no automorphism ambiguity.
Theorem~\ref{thm:ramsey-exact} applies.

We verify the amalgamation hypothesis explicitly. For a span of finite
ordered graphs, take the union over the common induced graph, put no new
edges between the two new parts, and extend the two agreeing linear orders
to a linear order of the union. This is an induced order-preserving amalgam.
Every triangle in this union lies in one of the original graphs, so the
construction also works in the triangle-free class. Disjoint ordered union
gives joint embedding. Both categories fail ordinary filteredness: the two
embeddings of a one-vertex graph into a fixed two-vertex graph remain
distinct after every injective map. Integral affine equalizers give affine
equalizers after every change of commutative coefficients. Totality follows
from the exactness implication in R2. The same argument proves the stated
general consequence for rigid Ramsey amalgamation classes.
\end{proof}

\begin{remark}[Unresolved: scope of the Ramsey argument]\label{C:ramsey-scope}
We have not proved that totality implies the finite-family Ramsey property.
Thus Theorem~\ref{thm:ramsey-exact} closes the named route C4 candidates,
but does not settle the main implication for arbitrary total categories.
We make no priority claim: searches for ``Ramsey exact colimits'',
``Ramsey affinization'', ``Ramsey Isbell Mitchell'', and related phrases
did not identify this implication, and Isbell's 1974 full text was not obtained.
\end{remark}
